\documentclass[11pt,letterpaper]{article}

\usepackage[top=1in, bottom=1in, left=1in, right=1in, headheight=14pt]{geometry}
\usepackage{fontspec}
\setmainfont{DejaVu Serif}
\setsansfont{DejaVu Sans}
\setmonofont{DejaVu Sans Mono}

\usepackage{xcolor}
\usepackage{titlesec}
\usepackage{fancyhdr}
\usepackage{enumitem}
\usepackage{hyperref}
\usepackage{booktabs}
\usepackage{tabularx}
\usepackage{longtable}
\usepackage{colortbl}
\usepackage{multirow}
\usepackage{array}
\usepackage{parskip}
\usepackage{tcolorbox}
\usepackage{graphicx}
\usepackage{lastpage}

\definecolor{primary}{HTML}{1B5E20}
\definecolor{secondary}{HTML}{1565C0}
\definecolor{tertiary}{HTML}{4E342E}
\definecolor{accent}{HTML}{2E7D32}
\definecolor{lightprimary}{HTML}{E8F5E9}
\definecolor{lightsecondary}{HTML}{E3F2FD}
\definecolor{lighttertiary}{HTML}{EFEBE9}
\definecolor{rowgray}{HTML}{F5F5F5}
\definecolor{bordergray}{HTML}{BDBDBD}
\definecolor{subtlegray}{HTML}{757575}
\definecolor{darktext}{HTML}{212121}

\titleformat{\section}
  {\Large\bfseries\sffamily\color{primary}}
  {\thesection}{1em}{}
  [\vspace{-0.5em}\textcolor{bordergray}{\rule{\textwidth}{0.4pt}}]

\titleformat{\subsection}
  {\large\bfseries\sffamily\color{secondary}}
  {\thesubsection}{1em}{}

\titleformat{\subsubsection}
  {\normalsize\bfseries\sffamily\color{tertiary}}
  {\thesubsubsection}{1em}{}

\titlespacing*{\section}{0pt}{1.5em}{0.8em}
\titlespacing*{\subsection}{0pt}{1.2em}{0.5em}
\titlespacing*{\subsubsection}{0pt}{0.8em}{0.3em}

\newcommand{\stageheader}[2]{%
  \noindent\colorbox{#2}{\makebox[\textwidth][l]{%
    \textcolor{white}{\bfseries\sffamily\hspace{6pt}#1\hspace{6pt}}}}%
  \vspace{0.5em}\par%
}

\newtcolorbox{asciibox}{
  colback=rowgray, colframe=bordergray,
  arc=1pt, boxrule=0.5pt,
  left=6pt, right=6pt, top=4pt, bottom=4pt,
  fontupper=\small\ttfamily,
  breakable
}

\newtcolorbox{quotebox}{
  colback=lightprimary, colframe=primary,
  arc=2pt, boxrule=0.5pt,
  left=12pt, right=12pt, top=8pt, bottom=8pt,
  breakable
}

\newcommand{\metafield}[2]{\textbf{\sffamily #1} #2\par}
\newcolumntype{L}[1]{>{\raggedright\arraybackslash}p{#1}}
\newcolumntype{C}[1]{>{\centering\arraybackslash}p{#1}}
\newcommand{\tableheader}[1]{\textbf{\sffamily\textcolor{white}{#1}}}

\pagestyle{fancy}
\fancyhf{}
\fancyhead[L]{\small\sffamily\textcolor{subtlegray}{PNW Fungi Taxonomy}}
\fancyhead[R]{\small\sffamily\textcolor{subtlegray}{GSD Mission Package}}
\fancyfoot[C]{\small\sffamily\textcolor{subtlegray}{Page \thepage\ of \pageref{LastPage}}}
\renewcommand{\headrulewidth}{0.4pt}
\renewcommand{\footrulewidth}{0pt}

\hypersetup{
  colorlinks=true,
  linkcolor=secondary,
  urlcolor=secondary,
  pdfauthor={GSD Ecosystem},
  pdftitle={Pacific Northwest Fungi Taxonomy -- Mission Package},
  pdfsubject={Vision to Mission Pipeline},
}

\begin{document}

\thispagestyle{empty}
\begin{center}
\vspace*{3cm}

{\small\sffamily\textcolor{primary}{\textbf{GSD RESEARCH MISSION PACKAGE}}}

\vspace{0.3cm}
\textcolor{bordergray}{\rule{8cm}{0.4pt}}

\vspace{1.5cm}
{\fontsize{28}{34}\selectfont\bfseries\sffamily\textcolor{primary}{The Quiet Light\\[0.3em]in the Dark Woods}}

\vspace{0.8cm}
{\Large\sffamily\textcolor{secondary}{Pacific Northwest Fungi Taxonomy}}

\vspace{0.5cm}
{\large\sffamily\itshape\textcolor{tertiary}{A Deep Research Study into the Fungal Kingdom of Cascadia}}

\vspace{3cm}
{\sffamily\textcolor{accent}{Vision $\rightarrow$ Research $\rightarrow$ Mission}}\\[0.3em]
{\small\sffamily\textcolor{accent}{Full Pipeline Execution}}

\vspace{3cm}
{\small\sffamily\textcolor{subtlegray}{Date: March 7, 2026 ~|~ Status: Ready for GSD Orchestrator}}\\[0.3em]
{\small\sffamily\textcolor{subtlegray}{Activation Profile: Squadron ~|~ Estimated: 14 context windows across 6 sessions}}
\end{center}
\newpage

\tableofcontents
\newpage

% =====================================================
% STAGE 1
% =====================================================
\stageheader{STAGE 1 --- VISION DOCUMENT}{primary}

\section{Pacific Northwest Fungi Taxonomy --- Vision Guide}

\metafield{Date:}{March 7, 2026}
\metafield{Status:}{Deep Research Vision / Pre-Mission}
\metafield{Depends on:}{gsd-foxfire-heritage-skills-vision.md, gsd-foxfire-heritage-mission-package.md, 04-salish-sea-expansion-vision.md, gsd-skill-creator-analysis.md, gsd-chipset-architecture-vision.md}
\metafield{Context:}{A comprehensive research mission exploring the taxonomy, ecology, ethnomycology, and conservation of fungi in the Pacific Northwest bioregion --- from the temperate rainforests of the Olympic Peninsula to the old-growth Douglas-fir stands of the Cascades, from the coastal spruce forests of British Columbia to the mixed-conifer woodlands of the Klamath Mountains. This pack treats mycological knowledge as both scientific discipline and heritage wisdom, bridging Western taxonomy with Indigenous ethnomycological traditions, and connecting citizen science with professional research through the Foxfire methodology of ``document what you find.''}

\subsection{Vision}

In 1966, when Eliot Wigginton's students in northeast Georgia named their magazine after the bioluminescent glow of fungi on decaying logs --- \textit{foxfire} --- they were naming something older than their mountains. The quiet light in the dark woods. The glow that appears on rotting wood after rain, produced by enzymes in the mycelium of \textit{Panellus stipticus} and \textit{Omphalotus olearius} and a handful of other species that evolution decided, for reasons still incompletely understood, should be luminous. The students probably did not know that more than 125 species of fungi produce bioluminescence. They chose the name because it was beautiful and strange and because it came from the forest floor, which is where the real knowledge always lives.

The Pacific Northwest is a mycological wonderland. The temperate rainforests of the Olympic Peninsula, the ancient Douglas-fir stands of the western Cascades, the fog-drenched Sitka spruce forests of the Oregon and Washington coasts, the cedar-hemlock forests of Vancouver Island --- these ecosystems harbor one of the densest concentrations of fungal diversity on Earth. The region's combination of maritime climate, volcanic soils, extreme precipitation gradients, and ancient forest continuity has produced a fungal flora of staggering richness: over 3,000 macrofungal species documented, with conservative estimates suggesting that fewer than one-third of the region's actual fungal species have been catalogued. The Pacific Northwest Fungi Project at Washington State University remains the only ongoing effort to systematically catalog the region's fungi --- there is no other state, federal, or private program doing this work.

But fungi are not merely species to be catalogued. Beneath every old-growth forest in the Pacific Northwest runs a mycorrhizal network --- what the German forester Peter Wohlleben called the ``wood-wide web'' --- through which trees exchange carbon, nitrogen, water, and chemical defense signals. Suzanne Simard's research on interior Douglas-fir forests in British Columbia has shown that these networks have a scale-free topology, with the largest, oldest trees serving as hubs connecting dozens of younger trees through shared fungal partners, primarily \textit{Rhizopogon vesiculosus} and \textit{R.~vinicolor}. The mycorrhizal network retains approximately 30\% of the sugars that connected trees produce through photosynthesis --- not parasitism, but payment for services rendered. When a tree is stressed, its neighbors send it nutrients through the network. When a tree is dying, it dumps its carbon reserves into the network for redistribution. The forest is not a collection of individual trees. It is a single organism connected by fungi, and the fungi are the connective tissue.

This is the space between. The literal, physical, biological space between the roots of one tree and the roots of another, bridged by mycelial threads one cell thick, running through soil that is itself largely composed of fungal biomass. Every principle that \textit{The Space Between} articulates mathematically --- that the relationships between things matter more than the things themselves, that networks create emergent properties their components cannot possess alone, that the spaces between define the structure --- every one of these principles is enacted daily, silently, beneath the forest floor of every old-growth stand in Cascadia.

\subsection{Problem Statement}

\begin{enumerate}[leftmargin=2em]
\item \textbf{The Taxonomic Deficit.} Globally, between 2.2 and 3.8 million fungal species are predicted to exist (Hawksworth \& Lücking 2018). Fewer than 150,000 have been formally described. In the Pacific Northwest specifically, fewer than one-third of predicted species have been catalogued. The Pacific Northwest Fungi Project is the sole ongoing effort --- no state, federal, or private program systematically inventories the region's fungi. Species are going extinct before they are named.

\item \textbf{The Identification Crisis.} Fungi have simple body plans with convergent morphological features, making visual identification unreliable for many groups. DNA barcoding using the ITS (Internal Transcribed Spacer) region has become the primary identification tool, but reference databases remain incomplete. Many environmental sequences cannot be matched to any known species. The trade-off between speed and accuracy in metabarcoding studies means that phylotype diversity is not necessarily a proxy for species richness.

\item \textbf{The Ethnomycological Erasure.} Indigenous Peoples of the Pacific Northwest have maintained relationships with fungi for millennia --- approximately 30--40 species are documented as having cultural roles for Canadian Indigenous groups alone, spanning food, medicine, fire-starting, mask-carving, and spiritual practice. Nancy J. Turner's collaborative research with Indigenous knowledge holders documents uses of fungi by Nlaka'pamux, Stl'atl'imx, Secwepemc, Nuu-chah-nulth, Kwakwaka'wakw, and many other nations. This knowledge is declining with elder loss and has not been systematically integrated with Western taxonomic frameworks.

\item \textbf{The Network Disruption.} Clearcutting, prescribed burning, and site preparation practices sever mycorrhizal networks that may have taken centuries to develop. USDA Forest Service research documents that disturbance reduces ectomycorrhizal formation and forest regeneration, with severity depending on disturbance type, mycorrhizal diversity, and site conditions. Invasion by non-ectomycorrhizal plants (e.g., \textit{Ceanothus}, \textit{Rubus}) can permanently alter soil microbial communities, making reforestation of Pinaceae species difficult or impossible on some sites.

\item \textbf{The Citizen Science Gap.} The Pacific Northwest has the densest concentration of amateur mycological societies in North America --- Puget Sound Mycological Society (50+ years), Oregon Mycological Society, Northwest Mushroomers Association, Cascade Mycological Society, and dozens more. These societies produce thousands of observations annually. But the Fungal Diversity Survey (FunDiS) has identified that much citizen science data is unusable because photographs are insufficient, identifications are incorrect, and critical substrate and habitat data are missing.
\end{enumerate}

\subsection{Core Concept}

\begin{quotebox}
\textbf{One-line model:} The Pacific Northwest fungal kingdom is a network of networks --- taxonomic, ecological, cultural, and scientific --- and this research mission documents each layer while teaching how to observe, identify, preserve, and connect them through the Foxfire methodology of community-based documentation.
\end{quotebox}

\subsubsection{Bioregion Map}

\begin{asciibox}
PACIFIC NORTHWEST FUNGI BIOREGION
===================================

  BRITISH COLUMBIA
  Vancouver Island     Coastal Ranges
  (cedar-hemlock)      (Sitka spruce)
       |                    |
  Strait of Georgia    Fraser Valley
       |                    |
  San Juan Islands     North Cascades
       |                    |
  OLYMPIC PENINSULA    WESTERN CASCADES
  (temperate rain-     (Douglas-fir /
   forest, Hoh/Quinault western hemlock)
   200+ in/yr precip)       |
       |               MT. RAINIER
  WILLAPA HILLS        subalpine meadows
  (coastal spruce)          |
       |               SOUTHERN CASCADES
  OREGON COAST RANGE   (mixed conifer)
  (Sitka spruce /           |
   western hemlock)    KLAMATH MOUNTAINS
       |               (mixed evergreen,
  WILLAMETTE VALLEY    Mediterranean/
  (oak woodlands,      maritime transition)
   prairie remnants)        |
       |               NORTHERN CALIFORNIA
  SISKIYOU MTNS        (Redwood Coast,
  (biodiversity         coastal fog belt)
   hotspot, serpentine)
\end{asciibox}

\subsubsection{Module Architecture}

\begin{asciibox}
MODULE ARCHITECTURE
===================

  [M1: TAXONOMY]          [M2: ECOLOGY]
  Classification          Mycorrhizal Networks
  Major Phyla/Orders      Saprotrophic Cycles
  PNW Key Species         Host Associations
  Morphological Keys      Forest Succession
        |                       |
        +--------.  .----------+
                 |  |
           [M3: ETHNOMYCOLOGY]
           Indigenous Knowledge
           Traditional Uses
           Cultural Sovereignty
           Nation-Specific
                 |
        +--------.  .----------+
        |                       |
  [M4: CONSERVATION]     [M5: APPLIED MYCOLOGY]
  Threatened Species      Mycoremediation
  Habitat Protection      Mycofiltration
  FunDiS Rare Fungi       Mycoforestry
  Climate Impacts         Cultivation
        |                       |
        +--------.  .----------+
                 |  |
         [M6: CITIZEN SCIENCE]
         DNA Barcoding Methods
         iNaturalist Protocols
         FunDiS Standards
         Foxfire Documentation
\end{asciibox}

\subsection{Research Modules}

\subsubsection{Module 1: Taxonomy \& Classification}

Systematic treatment of PNW macrofungi organized by phylum, order, and family. Covers Basidiomycota (agarics, boletes, polypores, coral fungi, tooth fungi, puffballs, jelly fungi) and Ascomycota (cup fungi, morels, truffles, flask fungi). Special attention to genera where PNW diversity is exceptional: \textit{Cortinarius}, \textit{Russula}, \textit{Ramaria}, \textit{Phaeocollybia}, \textit{Hygrocybe}, \textit{Tricholoma}, and the hypogeous genera \textit{Tuber}, \textit{Rhizopogon}, \textit{Leucangium}. Incorporates current molecular phylogenetic revisions.

\subsubsection{Module 2: Mycorrhizal Networks \& Forest Ecology}

The ecology of ectomycorrhizal and vesicular-arbuscular mycorrhizal associations in PNW forests. Documents the Simard research program on common mycorrhizal networks in Douglas-fir forests, the Beiler topology studies showing scale-free network architecture, the Molina/Trappe work on host specificity, and the Amaranthus research on disturbance effects. Covers the truffle--small-mammal--spotted-owl food web and defense signal transmission through mycelial networks.

\subsubsection{Module 3: Ethnomycology \& Traditional Knowledge}

Indigenous relationships with fungi across the PNW bioregion, organized by nation. Drawing on Turner's collaborative research with BC First Nations, Anderson \& Lake's California Indian ethnomycology, and the Fungi Foundation's Ethnomycology Ethical Guidelines (2025). Documents food, material, medicinal, and spiritual uses. All content governed by OCAP/CARE/UNDRIP frameworks.

\subsubsection{Module 4: Conservation \& Threatened Species}

Conservation status of PNW fungi including the FunDiS West Coast Rare Fungi Challenge, IUCN assessments, and the movement to recognize fungi as their own kingdom deserving independent conservation. Documents habitat threats and the Agarikon (\textit{Fomitopsis officinalis}) as an old-growth flagship species.

\subsubsection{Module 5: Applied Mycology}

Mycoremediation (Stamets' petroleum hydrocarbon experiments), mycofiltration (MycoBooms for Puget Sound), mycoforestry (post-disturbance recovery), and mycopesticides. Connects to heritage skills through accessible cultivation techniques.

\subsubsection{Module 6: Citizen Science \& Modern Methods}

DNA barcoding methodology, iNaturalist protocols, FunDiS quality standards, and Foxfire documentation applied to mycological field work. Covers miniPCR educational kits, MycoMatch AI identification, and the North American Mycoflora vision.

\subsection{Chipset Configuration}

\begin{asciibox}
mission:
  name: pnw-fungi-taxonomy
  version: 1.0
  activation: squadron
  crew_size: 12
  parallel_tracks: 2

models:
  opus: 30%   # Synthesis, ethnomycology, integration
  sonnet: 60% # Surveys, taxonomy, ecology, publication
  haiku: 10%  # Schema scaffolding, citation templates

safety:
  warden: true
  cultural_sovereignty: true
  species_location: redact_precise
  identification_disclaimer: mandatory
\end{asciibox}

\subsection{Scope Boundaries}

\subsubsection{In Scope (v1.0)}

\begin{itemize}[leftmargin=2em]
\item Macrofungi (visible fruiting bodies) of the PNW bioregion: BC coast through northern California
\item Ectomycorrhizal and VA mycorrhizal ecology in PNW forest ecosystems
\item Ethnomycological knowledge documented through published collaborative research
\item DNA barcoding methodology and citizen science protocols
\item Mycoremediation, mycofiltration, and mycoforestry applications
\item Conservation status assessment and rare species documentation
\item Integration with existing Heritage Skills Educational Pack architecture
\end{itemize}

\subsubsection{Out of Scope}

\begin{itemize}[leftmargin=2em]
\item Microfungi (molds, yeasts, rusts, smuts) except where ecologically significant
\item Psychoactive species identification or cultivation guidance
\item Medical claims beyond published peer-reviewed research
\item Precise GPS coordinates of rare or threatened species
\item Unpublished or non-consensual Indigenous knowledge
\item Commercial foraging guidance or harvest regulations
\end{itemize}

\subsection{Success Criteria}

\begin{enumerate}[leftmargin=2em]
\item Taxonomy module covers $\geq$100 key PNW macrofungal species with current nomenclature
\item Ecology module documents $\geq$5 major mycorrhizal network studies with quantitative findings
\item Ethnomycology module references $\geq$8 specific nations by name with documented fungal relationships
\item All ethnomycological content passes Cultural Sovereignty Warden review
\item Conservation module includes the complete FunDiS West Coast Rare Fungi Challenge species list
\item Applied mycology module documents $\geq$3 mycoremediation case studies with measured outcomes
\item Citizen science module provides complete DNA barcoding workflow from collection to sequence submission
\item Safety Warden enforces mandatory identification disclaimers on all edibility references
\item Cross-module synthesis demonstrates $\geq$4 explicit connections
\item All numerical claims attributed to specific peer-reviewed or government sources
\item Package compiles and executes in GSD skill-creator with zero manual intervention
\item Through-line connects mycorrhizal networks to ``the spaces between'' ecosystem philosophy
\end{enumerate}

\subsection{Relationship to Other Documents}

\begin{center}
\rowcolors{2}{rowgray}{white}
\begin{tabularx}{\textwidth}{L{4.5cm}X}
\toprule
\rowcolor{primary}
\tableheader{Document} & \tableheader{Relationship} \\
\midrule
gsd-foxfire-heritage-skills-vision.md & Parent heritage pack; ``Foxfire'' derives from bioluminescent fungi \\
04-salish-sea-expansion-vision.md & Cedar Path badges require cedar--fungus symbiosis understanding \\
gsd-foxfire-heritage-mission-package.md & SUMO ontology extension for mycological knowledge \\
gsd-fair-use-educational-mission.md & Governs use of canonical mycological works \\
The Space Between (Tibsfox) & Mycorrhizal networks as scale-free graphs; fungal growth as L-systems \\
\bottomrule
\end{tabularx}
\end{center}

\subsection{The Through-Line}

Fungi are the original space between.

Long before anyone wrote about the relationships between things mattering more than the things themselves, fungi were enacting this principle underground. A mycorrhizal network is not a metaphor for connection --- it is connection, physical and measurable, carbon atoms moving from one organism to another through hyphal threads one cell thick, across distances that matter ecologically, between organisms that could not survive alone. The mycorrhizal network retains its 30\% fee not as parasitism but as infrastructure maintenance --- the cost of keeping the space between functional.

The Foxfire students in Georgia named their magazine after the glow of fungi on rotting logs without knowing they were naming the connective tissue of the forest. This research mission honors that naming by documenting the full scope of what fungi do, what they connect, what they mean to the peoples who have lived alongside them for millennia, and what they can teach us about building systems --- technological, social, ecological --- that are resilient because they are connected, not because they are strong.

\newpage

% =====================================================
% STAGE 2
% =====================================================
\stageheader{STAGE 2 --- RESEARCH REFERENCE}{secondary}

\section{Pacific Northwest Fungi --- Research Reference}

\subsection{How to Use This Document}

This reference compiles findings from peer-reviewed research, government agencies, professional mycological organizations, and published collaborative ethnomycological studies. Key source organizations: USDA Forest Service PNW Research Station; Pacific Northwest Fungi Project (WSU); Puget Sound Mycological Society; Fungal Diversity Survey (FunDiS); Fungi Perfecti; UNITE Database; iNaturalist; Mycological Society of America; Nancy J. Turner (University of Victoria).

\subsection{Module 1 Reference: Taxonomy \& Classification}

\subsubsection{Scope of PNW Fungal Diversity}

The Pacific Northwest harbors exceptional macrofungal diversity. \textit{Mushrooms of Cascadia} by Siegel and Schwarz (Backcountry Press, 576 pp.) treats over 750 species. Trudell and Ammirati's \textit{Mushrooms of the Pacific Northwest} covers 460 species. Globally, approximately 3 million fungal species are predicted; only about 150,000 have been described. In the PNW, fewer than one-third of predicted species have been catalogued (Pacific Northwest Fungi Project, WSU).

\subsubsection{Major Taxonomic Groups}

\begin{center}
\rowcolors{2}{rowgray}{white}
\begin{tabularx}{\textwidth}{L{2.5cm}L{3.5cm}X}
\toprule
\rowcolor{primary}
\tableheader{Group} & \tableheader{Key PNW Genera} & \tableheader{Ecological Role} \\
\midrule
Agarics & \textit{Cortinarius, Russula, Amanita, Tricholoma, Phaeocollybia} & Ectomycorrhizal partners of conifers \\
Boletes & \textit{Boletus, Suillus, Leccinum} & Ectomycorrhizal; \textit{B.~edulis} commercially harvested \\
Chanterelles & \textit{Cantharellus, Craterellus} & Ectomycorrhizal; \textit{C.~formosus} is Oregon's state mushroom \\
Polypores & \textit{Fomitopsis, Ganoderma, Trametes} & Saprotrophic; \textit{F.~officinalis} old-growth indicator \\
Coral fungi & \textit{Ramaria, Clavulina} & Mixed; PNW is a center of \textit{Ramaria} diversity \\
Truffles & \textit{Tuber, Rhizopogon, Leucangium} & Ectomycorrhizal; critical small mammal food \\
Morels & \textit{Morchella} & Post-fire fruiting specialists \\
\bottomrule
\end{tabularx}
\end{center}

\subsubsection{DNA Barcoding}

The ITS region (~600 bp) is the universal fungal DNA barcode (Schoch et al. 2012). The UNITE database provides curated sequences. Predicted similarity cutoffs vary significantly between clades, so no single threshold reliably delineates species. Nanopore long-read sequencing now enables ITS--LSU amplicons up to 2,758 bp for higher resolution. Integrative taxonomy combining molecular, morphological, and ecological data remains the gold standard (Lücking et al. 2020).

\subsection{Module 2 Reference: Mycorrhizal Networks}

\subsubsection{Ectomycorrhizal Associations}

Most PNW forest trees require ectomycorrhizae for nutrient and water uptake (USDA Forest Service). Major host families: Pinaceae, Fagaceae, Betulaceae, Salicaceae. VAM fungi formed by maples, cedars, and shrubs are incompatible with Pinaceae --- sites dominated by VAM species gradually lose the ectomycorrhizal fungi conifers require.

\subsubsection{Common Mycorrhizal Networks}

Simard et al.'s research documents that most trees in uneven-aged Douglas-fir forests are interconnected by networks of \textit{Rhizopogon vesiculosus} and \textit{R.~vinicolor}. These networks exhibit scale-free topology with hub trees maintaining the majority of connections. Carbon transfer flows from canopy trees to shaded seedlings. The network retains approximately 30\% of photosynthetic sugars as mutualistic payment.

\subsubsection{Truffle--Mammal--Owl Food Web}

Oregon white truffle (\textit{Tuber gibbosum}) associates exclusively with Douglas-fir. Truffles are critical food for squirrels, voles, and chipmunks, which are major prey for the northern spotted owl. This three-trophic-level dependency means mycorrhizal disruption cascades to vertebrate conservation targets.

\subsection{Module 3 Reference: Ethnomycology}

Turner's collaborative research documents approximately 30--40 fungal species with cultural roles for Canadian Indigenous Peoples. Food uses include chanterelles, matsutake, morels, and puffballs. Material uses include \textit{Fomes fomentarius} for fire-starting and \textit{Ganoderma applanatum} for mask carving. Puffball spores are widely documented for treating wounds. Anderson \& Lake document California Indian tribes including Karuk, Hupa, Yurok, and Nisenan Maidu utilizing fungi, with the Karuk deliberately spreading matsutake spores to extend fruiting areas.

\subsection{Module 4 Reference: Conservation}

FunDiS West Coast Rare Fungi Challenge targets 20 rare or under-documented species. \textit{Fomitopsis officinalis} (Agarikon) is an old-growth flagship. \textit{Pleurotus citrinopileatus} (Golden Oyster) has escaped cultivation and threatens native communities. Climate projections suggest increased drought and disturbance threatening ectomycorrhizal communities adapted to mesic old-growth conditions.

\subsection{Module 5 Reference: Applied Mycology}

Stamets' WSDOT experiment: petroleum hydrocarbons at ~20,000 ppm reduced below 200 ppm using \textit{Pleurotus ostreatus}. MycoBooms tested in Puget Sound for estuary mycofiltration. Over 120 novel enzymes identified from mushroom-forming fungi. Mycoforestry: retention of ectomycorrhizal hosts after disturbance is critical for reforestation (Amaranthus \& Perry, USDA Forest Service).

\subsection{Module 6 Reference: Citizen Science}

PNW hosts the densest network of amateur mycological societies in North America. FunDiS has sequenced approximately 5,000 specimens, discovering undescribed species in \textit{Amanita}, \textit{Russula}, and \textit{Callistosporium}. The Ya'lis Alert Bay survey exemplifies place-based citizen mycology on Kwakwaka'wakw territory. Complete barcoding workflow: collection $\rightarrow$ voucher $\rightarrow$ tissue preservation $\rightarrow$ DNA extraction $\rightarrow$ ITS PCR $\rightarrow$ sequencing $\rightarrow$ database submission.

\subsection{Safety Considerations}

\begin{center}
\rowcolors{2}{rowgray}{white}
\begin{tabularx}{\textwidth}{L{3cm}L{2cm}X}
\toprule
\rowcolor{primary}
\tableheader{Area} & \tableheader{Boundary} & \tableheader{Implementation} \\
\midrule
Edibility claims & ABSOLUTE & Mandatory identification disclaimer \\
Toxic look-alikes & ABSOLUTE & Every edible entry includes toxic mimics \\
Rare species locations & ABSOLUTE & No precise coordinates \\
Indigenous knowledge & GATE & Published collaborative research only \\
Medical claims & GATE & Peer-reviewed sources only \\
Psychoactive species & ANNOTATE & Acknowledge existence; no ID keys \\
\bottomrule
\end{tabularx}
\end{center}

\subsection{Source Bibliography}

\subsubsection{Government \& Agency Sources}
Molina, R. et al. 1994. USDA Forest Service Gen. Tech. Rep. PSW-151.
Amaranthus, M.P. 1992. USDA Forest Service PNW Research Station.
Pilz, D. \& Molina, R. 2002. Forest Ecology and Management 155: 3--16.

\subsubsection{Peer-Reviewed Research}
Beiler, K.J. et al. 2015. \textit{J. Ecol.} 103: 616--628.
Birch, J.D. et al. 2021. \textit{J. Ecol.} 109: 806--818.
Lücking, R. et al. 2020. \textit{IMA Fungus} 11:14.
Turner, N.J. 2021. \textit{Canadian Journal of Botany.}
Anderson, M.K. \& Lake, F.K. 2013. \textit{J. Ethnobiology} 33(1): 33--85.
Simard, S.W. 2009. \textit{Forest Ecology and Management} 258: 95--107.
Hawksworth, D.L. \& Lücking, R. 2018. \textit{Microbiology Spectrum} 5(4).

\subsubsection{Key Texts}
Siegel, N. \& Schwarz, C. 2023. \textit{Mushrooms of Cascadia.} Backcountry Press.
Trudell, S. \& Ammirati, J. 2009. \textit{Mushrooms of the Pacific Northwest.} Timber Press.
Stamets, P. 2005. \textit{Mycelium Running.} Ten Speed Press.
Arora, D. 1986. \textit{Mushrooms Demystified.} Ten Speed Press.
Fungi Foundation. 2025. \textit{Ethnomycology Ethical Guidelines.}

\newpage

% =====================================================
% STAGE 3
% =====================================================
\stageheader{STAGE 3 --- MISSION PACKAGE}{tertiary}

\section{Milestone Specification}

\subsection{Mission Objective}

Build the PNW Fungi Taxonomy research module as a complete skill-creator component: six interconnected modules, Safety Warden with identification disclaimer and cultural sovereignty rules, and integration hooks to the Heritage Skills Educational Pack. Serves as both standalone mycological reference and template for place-based natural history documentation.

\subsection{Architecture Overview}

\begin{asciibox}
WAVE EXECUTION ARCHITECTURE
============================

WAVE 0 (Sequential) -- Foundation
  [0.1] Shared Schema + Source Index -----> cache

WAVE 1 (Parallel -- 2 tracks)
  Track A:                    Track B:
  [1.1] Taxonomy Module       [1.2] Ecology Module
  [1.3] Conservation Module   [1.4] Ethnomycology Module

WAVE 2 (Parallel -- 2 tracks)
  Track A:                    Track B:
  [2.1] Applied Mycology      [2.2] Citizen Science

WAVE 3 (Sequential) -- Synthesis
  [3.1] Cross-Module Integration + Safety Warden

WAVE 4 (Sequential) -- Publication
  [4.1] Document Assembly + Verification
\end{asciibox}

\subsection{Deliverables}

\begin{center}
\rowcolors{2}{rowgray}{white}
\begin{tabularx}{\textwidth}{C{0.7cm}L{3cm}XL{1.5cm}}
\toprule
\rowcolor{primary}
\tableheader{\#} & \tableheader{Deliverable} & \tableheader{Criteria} & \tableheader{Spec} \\
\midrule
D1 & Taxonomy Reference & $\geq$100 species, current names & 1.1 \\
D2 & Ecology Reference & $\geq$5 CMN studies & 1.2 \\
D3 & Ethnomycology Ref. & $\geq$8 nations, OCAP & 1.4 \\
D4 & Conservation Ref. & FunDiS species list & 1.3 \\
D5 & Applied Mycology & $\geq$3 case studies & 2.1 \\
D6 & Citizen Science Guide & Full barcoding workflow & 2.2 \\
D7 & Safety Warden Rules & Disclaimers + cultural gates & 3.1 \\
D8 & Cross-Module Synthesis & $\geq$4 connections & 3.1 \\
D9 & Published Document & Compiled, indexed & 4.1 \\
\bottomrule
\end{tabularx}
\end{center}

\subsection{Component Breakdown}

\begin{center}
\rowcolors{2}{rowgray}{white}
\begin{tabularx}{\textwidth}{L{3cm}L{2cm}C{1.5cm}C{1.5cm}}
\toprule
\rowcolor{primary}
\tableheader{Component} & \tableheader{Deps} & \tableheader{Model} & \tableheader{Tokens} \\
\midrule
0.1 Schema + Index & None & Haiku & 4K \\
1.1 Taxonomy & 0.1 & Sonnet & 25K \\
1.2 Ecology & 0.1 & Sonnet & 20K \\
1.3 Conservation & 0.1 & Sonnet & 15K \\
1.4 Ethnomycology & 0.1 & Opus & 20K \\
2.1 Applied Mycology & 0.1 & Sonnet & 15K \\
2.2 Citizen Science & 0.1, 1.1 & Sonnet & 15K \\
3.1 Integration & All above & Opus & 20K \\
4.1 Publication & 3.1 & Sonnet & 10K \\
\bottomrule
\end{tabularx}
\end{center}

\section{Wave Execution Plan}

\begin{center}
\rowcolors{2}{rowgray}{white}
\begin{tabularx}{\textwidth}{C{1.2cm}L{3cm}C{1.5cm}C{2cm}C{2cm}}
\toprule
\rowcolor{primary}
\tableheader{Wave} & \tableheader{Purpose} & \tableheader{Tracks} & \tableheader{Components} & \tableheader{Windows} \\
\midrule
0 & Foundation & 1 (seq) & 0.1 & 1 \\
1 & Core Research & 2 (par) & 1.1--1.4 & 4 \\
2 & Applied Research & 2 (par) & 2.1--2.2 & 2 \\
3 & Synthesis & 1 (seq) & 3.1 & 2 \\
4 & Publication & 1 (seq) & 4.1 & 1 \\
\bottomrule
\end{tabularx}
\end{center}

\subsection{Dependency Graph}

\begin{asciibox}
[0.1 Schema] --+----+----+----+----+
               |    |    |    |    |
               v    v    v    v    v
           [1.1] [1.2] [1.3] [1.4] |
               |    |    |    |    |
               +----+----+----+    v
               |              [2.1 App]
               v
           [2.2 CitSci]
               |
       All ----+----> [3.1 Integration]
                           |
                           v
                     [4.1 Publication]
\end{asciibox}

\section{Test Plan}

\subsection{Safety-Critical Tests}

\begin{center}
\rowcolors{2}{rowgray}{white}
\begin{tabularx}{\textwidth}{C{0.8cm}L{3cm}X}
\toprule
\rowcolor{primary}
\tableheader{ID} & \tableheader{Verifies} & \tableheader{Expected Result} \\
\midrule
S1 & Edibility disclaimer & Every edibility mention has warning \\
S2 & Toxic look-alikes & Every edible entry includes mimics \\
S3 & Rare species redaction & No precise coordinates \\
S4 & Cultural sovereignty & Nation-specific names; published sources \\
S5 & Medical claim sourcing & All claims peer-reviewed \\
S6 & Psychoactive boundary & No ID keys or cultivation guidance \\
\bottomrule
\end{tabularx}
\end{center}

\subsection{Verification Matrix}

\begin{center}
\rowcolors{2}{rowgray}{white}
\begin{tabularx}{\textwidth}{C{0.8cm}XL{2cm}C{1.2cm}}
\toprule
\rowcolor{primary}
\tableheader{\#} & \tableheader{Criterion} & \tableheader{Tests} & \tableheader{Status} \\
\midrule
1 & $\geq$100 species & C1, C2 & -- \\
2 & $\geq$5 CMN studies & C3, C4 & -- \\
3 & $\geq$8 nations & C5, S4 & -- \\
4 & OCAP compliance & S4, I1 & -- \\
5 & FunDiS list & C6 & -- \\
6 & $\geq$3 remediation cases & C7 & -- \\
7 & Barcoding workflow & C9 & -- \\
8 & Safety disclaimers & S1, S2, S5 & -- \\
9 & $\geq$4 connections & I2, I3 & -- \\
10 & Claims sourced & C11, S5 & -- \\
11 & GSD executable & I4, I5 & -- \\
12 & Through-line & I6 & -- \\
\bottomrule
\end{tabularx}
\end{center}

\section{Crew Manifest}

\begin{center}
\rowcolors{2}{rowgray}{white}
\begin{tabularx}{\textwidth}{L{2.5cm}L{2.5cm}X}
\toprule
\rowcolor{primary}
\tableheader{Role} & \tableheader{Agent} & \tableheader{Responsibility} \\
\midrule
Commander & MYCELIA-CMD & Mission coordination \\
Chief Scientist & TAXONOMY-LEAD & Taxonomic accuracy \\
Ecology Lead & NETWORK-LEAD & Mycorrhizal research \\
Cultural Liaison & ETHNO-LEAD & Ethnomycology, sovereignty \\
Conservation & CONSERVE-LEAD & Threat assessment \\
Applied Research & APPLIED-LEAD & Mycoremediation cases \\
Methods Lead & CITIZEN-LEAD & Barcoding, protocols \\
Safety Warden & SAFETY-WARDEN & Disclaimers, boundaries \\
Integration & SYNTHESIS-INT & Cross-module connections \\
QA & QA-VERIFY & Test execution \\
Doc Assembly & DOC-ASSEMBLE & LaTeX, formatting \\
Source Librarian & SOURCE-LIB & Citation validation \\
\bottomrule
\end{tabularx}
\end{center}

\section{Execution Summary}

\begin{center}
\rowcolors{2}{rowgray}{white}
\begin{tabularx}{\textwidth}{L{4.5cm}X}
\toprule
\rowcolor{primary}
\tableheader{Metric} & \tableheader{Value} \\
\midrule
Total components & 9 \\
Parallel tracks (max) & 2 \\
Sequential depth & 4 waves \\
Context windows & 14 \\
Sessions & 6 \\
Token budget & $\sim$144K \\
Opus allocation & 28\% \\
Sonnet allocation & 69\% \\
Haiku allocation & 3\% \\
Safety-critical tests & 6 (all BLOCK) \\
Activation profile & Squadron \\
\bottomrule
\end{tabularx}
\end{center}

\vspace{1em}
\begin{quotebox}
\textit{``The students chose to call their magazine Foxfire --- after the bioluminescent glow of fungi on decaying logs in the North Georgia forests. A quiet light in the dark woods. They did not know that they were naming the connective tissue of the forest itself --- the mycelium that runs through every rotting log, every handful of forest soil, every space between one living thing and another. The light was always there. They just noticed it.''}

\vspace{0.5em}
\hfill --- The Through-Line
\end{quotebox}

\end{document}
